% Autoría principal: Mario Zambrano.

\subsection{Clasificación de fallos}

En un sistema distribuido, los componentes pueden
presentar fallos de distintas formas. Un servidor
puede dejar de funcionar, perder una comunicación
o responder fuera del tiempo esperado. Por eso,
antes de establecer una solución, es importante
identificar qué tipo de fallo se quiere tratar.
Nogueira y Coelho \cite{nogueira2026fault}
explican que las técnicas de recuperación dependen
del modelo de fallos considerado.

Entre los principales tipos se encuentran:

\begin{itemize}

    \item \textbf{Fallo por caída (\textit{crash}):}
    ocurre cuando un proceso deja de ejecutarse
    y ya no responde a las solicitudes. En nuestra
    práctica, una situación de este tipo sería
    detener el servidor del nodo-B mientras otro
    nodo intenta comunicarse con él.

    \item \textbf{Fallo por omisión:}
    se presenta cuando un proceso no realiza un
    envío o una recepción que debía efectuar.
    Para el cliente, uno de los posibles síntomas
    es no recibir la respuesta esperada, aunque
    esto no permite determinar por sí solo la
    causa del problema.

    \item \textbf{Fallo de temporización:}
    aparece cuando una operación se completa
    fuera del intervalo de tiempo establecido.
    La respuesta puede ser correcta, pero llegar
    demasiado tarde para resultar útil.

    \item \textbf{Fallo bizantino:}
    se produce cuando un nodo presenta un
    comportamiento arbitrario. Por ejemplo,
    podría enviar información contradictoria
    a distintos participantes sin dejar de
    funcionar.
\end{itemize}

Los fallos por caída, omisión y temporización
forman parte de los modelos que Nogueira y
Coelho \cite{nogueira2026fault} distinguen
en su estudio sobre tolerancia a fallos.
Por su parte, Zhong et al.
\cite{zhong2023byzantine} explican el
problema de los nodos bizantinos al
analizar algoritmos que deben alcanzar
consenso incluso cuando existen
participantes defectuosos.

Nuestro prototipo permite observar algunos
errores de comunicación, pero no implementa
consenso bizantino ni mecanismos para
tolerar cualquier tipo de fallo.

\subsection{Detección, recuperación y fiabilidad}

Cuando un cliente no recibe una respuesta,
no siempre es posible saber qué ocurrió.
El servidor pudo detenerse, la comunicación
pudo interrumpirse o la solicitud pudo
tardar más de lo esperado. Duarte et al.
\cite{duarte2023detectors} explican que,
en los sistemas asíncronos, resulta difícil
distinguir entre un proceso que ha fallado
y otro que simplemente responde con lentitud.

Una forma de controlar estas situaciones es
establecer tiempos de espera. En nuestra
implementación, los clientes TCP y gRPC
utilizan un límite de dos segundos.
El cliente TCP captura excepciones
relacionadas con la conexión y la lectura
de mensajes, mientras que gRPC permite
identificar errores producidos durante
las llamadas remotas.

Estos errores pueden registrarse en los
archivos de latencia, lo que permite
separar las solicitudes correctas de
aquellas que terminaron con problemas.
Sin embargo, identificar un error no
significa que el servicio se haya
recuperado.

Entre las técnicas de recuperación se
encuentran los reintentos, el reinicio
de componentes y el uso de réplicas.
Nogueira y Coelho \cite{nogueira2026fault}
distinguen la recuperación de un servicio
de la recuperación de su estado, ya que
reiniciar un proceso no garantiza que
se conserven los datos o mensajes
pendientes.

En el trabajo de Marletta et al.
\cite{marletta2024robustness} se presenta
una solución que modifica llamadas remotas
en aplicaciones Java e incorpora políticas
de reintentos. Los autores muestran una
configuración con esperas progresivas y
un número máximo de intentos para
controlar fallos temporales.

En nuestro caso no se implementó una
política propia de reintentos automáticos
ni un mecanismo de recuperación mediante
réplicas. Los clientes pueden detectar y
registrar determinados errores, pero
todavía no cuentan con una solución que
garantice la continuidad del servicio
cuando un nodo deja de funcionar.

\subsection{Seguridad: JWT y OAuth 2.0}

La seguridad de un sistema distribuido
no depende únicamente de que los mensajes
lleguen correctamente. También es necesario
comprobar quién solicita una operación
y si tiene permiso para realizarla.

La autenticación permite comprobar la
identidad de un usuario o proceso,
mientras que la autorización determina
qué recursos puede utilizar. Bucko et al.
\cite{bucko2023jwt} explican esta
diferencia y presentan una aplicación
web que utiliza JWT como parte de su
mecanismo de control de acceso.

JSON Web Token (JWT) permite representar
información mediante un token que contiene
declaraciones sobre su portador. En las
implementaciones con firma, esta puede
utilizarse para comprobar que el contenido
no fue modificado. El token también puede
contener datos como el identificador del
usuario o su tiempo de expiración.
Bucko et al. \cite{bucko2023jwt}
describen su utilización en el proceso
de autenticación y autorización de
una aplicación web.

OAuth 2.0 cumple una función diferente.
Es un marco de autorización que permite
conceder acceso limitado a determinados
recursos sin entregar directamente las
credenciales del usuario a la aplicación
que solicita el acceso.

Almeida y Canedo
\cite{almeida2022authmicroservices}
identifican OAuth 2.0 y JWT entre los
mecanismos utilizados para proteger
aplicaciones basadas en microservicios.
También explican que ambos pueden
emplearse conjuntamente, aunque uno
establece un mecanismo de autorización
y el otro define un formato de token.

La seguridad de los tokens requiere
validaciones adecuadas y protección
durante su transmisión. Por ejemplo,
una firma digital no reemplaza el
cifrado del canal de comunicación.

En nuestra práctica, los mensajes TCP
se intercambian mediante JSON y las
llamadas gRPC utilizan Protocol Buffers.
No se implementaron JWT, OAuth 2.0
ni un mecanismo de autenticación de
usuarios. Además, las conexiones se
realizaron en localhost sin configurar
TLS.

Por tanto, estas tecnologías se estudian
como alternativas para proteger una
aplicación distribuida que deba gestionar
usuarios, permisos y acceso a servicios,
pero no forman parte de las funciones
implementadas en esta práctica.

\subsection{CAP y ejemplos de aplicaciones reales}

El teorema CAP explica una limitación
de los sistemas distribuidos que
mantienen información compartida entre
varios nodos. Sus tres propiedades son
consistencia, disponibilidad y tolerancia
a particiones de red.

La consistencia está relacionada con
mantener una visión coherente de los
datos compartidos. La disponibilidad
se refiere a la capacidad de responder
a las solicitudes, mientras que una
partición ocurre cuando se interrumpe
la comunicación entre grupos de nodos.

Durante una partición, un sistema no
puede garantizar al mismo tiempo
consistencia fuerte y disponibilidad
para todas las solicitudes. Lee et al.
\cite{lee2023cap} estudian esta
relación y muestran cómo los requisitos
de consistencia y disponibilidad
también se relacionan con la latencia
de comunicación.

Este problema puede observarse, por
ejemplo, en una aplicación que mantiene
copias del inventario de una tienda
en distintos servidores. Si estos
pierden la comunicación, aceptar
nuevas compras sin confirmar las
existencias podría ocasionar datos
inconsistentes. Una alternativa
sería detener temporalmente ciertas
operaciones hasta recuperar la
comunicación.

En cambio, una aplicación de
publicaciones podría permitir que
los usuarios continúen registrando
contenido durante una interrupción.
Para hacerlo tendría que aceptar
diferencias temporales entre las
copias y resolverlas posteriormente.

Los casos anteriores son ejemplos
hipotéticos. Su finalidad es mostrar
que las decisiones sobre consistencia
y disponibilidad dependen de las
necesidades de cada aplicación.

Es importante aclarar que CAP no
determina si conviene utilizar sockets
TCP o gRPC. Ambos son mecanismos de
comunicación, mientras que CAP trata
sobre las garantías de un sistema que
mantiene datos distribuidos.

En nuestra práctica se compararon
latencias y rendimiento de TCP y gRPC,
y se incorporaron relojes de Lamport
para registrar el orden de los eventos.
No se implementó una base de datos
replicada ni se evaluaron formalmente
las propiedades del teorema CAP.
Por esta razón, su aplicación en
el informe es conceptual.
